\subsection{Raggi cosmici e produzione dei muoni}
I raggi cosmici sono fondamentalmente nuclei accelerati da sorgenti astrofisiche. In particolare la parte bassa dello spettro dei raggi cosmici è dominata da nuclei di origine solare (vento solare). Nelle interazioni che avvengono nella parte alta dell'atmosfera vengono prodotte numerose particelle secondarie. Queste particelle secondarie possono a loro volta produrre altre particelle provocando uno sciame. Una parte di questa radiazione prodotta dal decadimento di pioni carichi si trasformano in muoni che possono essere rivelati sulla superficie terrestre.

I pioni possono essere positivi $\pi^+$, negativi $\pi^-$ o neutri $\pi^0$. Il pione neutro decade elettromagneticamente in due fotoni con un $\text{B.R.}=98.8\%$:
\begin{equation}
    \pi^0\,\to\,\gamma\gamma
\end{equation}
O con meno frequenza ($\text{B.R.}=1.2\%$) in fotone, elettrone e positrone:
\begin{equation}
    \pi^0\,\to\,\gamma\,+\,e^-\,+\,e^+
\end{equation}
I pioni carichi invece decadono in muoni e neutrini:
\begin{equation}
    \pi^+\,\to\,\mu^+\,+\,\nu_\mu \qquad \pi^-\,\to\,\mu^-\,+\,\bar{\nu}_\mu
\end{equation}
I decadimenti di interesse sono quelli dei pioni carichi. Una caratteristica fondamentale dei pioni è che hanno spin zero. Questa caratteristica, insieme al fatto che i neutrini hanno una ellicità ben definita (negativa per i neutrini e positiva per gli anti-neutrini),  a terra vedremo una polarizzazione dei muoni.
\subsection{I Muoni e la perdita di energia}
Il muone è una particella elementare con una massa di circa $m_\mu\approx105\,\unit{MeV}$ e spin $1/2$. Esistono muoni carichi positivamente $\mu^+$ e muoni negativi $\mu^-$. Il muone è un leptone instabile in quanto decade debolmente,  con una vita media di $\tau_\mu\approx2.2\,\mu s$. 
La legge di decadimento è:
\begin{equation}
    N(t)=N_0e^{-\frac{t}{\tau_\mu}}
\end{equation}
Dove $N(t)$ rappresenta il numero di particelle non ancora decadute, $N_0$ il numero di particelle iniziale e $\tau$ la vita media delle particelle.
Ad esempio, il muone carico positivamente, decade in un positrone, un neutrino e un anti-neutrino.
\begin{equation}
    \mu^+\,\,\to\,\,\nu_{e}\,+\,\bar{\nu}_e\,+\,e^+
\end{equation}
Essendo un decadimento a tre corpi, l'energia del positrone ha un end-point che corrisponde alla situazione in cui neutrino e anti-neutrino vengono emessi nello stesso emisfero mentre il positrone in quello opposto. Questa è una configurazione altamente probabile.

I muoni viaggiano ad una velocità prossima a quella della luce attraversando diverse decine di chilometri (a causa della dilatazione temporale) prima di decadere in $e,\nu,\bar{\nu}$. Per osservare questo decadimento occorre quindi fermarli. La perdita di energia è costante finchè la particella si muove con velocità prossima a quella della luce. Appena la velocità diminuisce , la particella subisce una brusca decelerazione (picco di Bragg). Una volta fermato nel materiale il muone negativo e il muone positivo decadono in tempi differenti. 

I muoni negativi, come i muoni positivi, perdono energia per interazioni multiple ma a causa della carica negativa una volta fermi risentono dell'attrazione dei nuclei degli atomi. Essendo particelle diverse non risentono del principio di esclusione di Pauli e vengono attratti verso il livello fondamentale dell'atomo emettendo raggi x. La vita media dei muoni negativi fermati nella materia è minore della vita media nel vuoto in quanto si aggiunge un altro effetto di "scomparsa".
Per quanto riguarda i muoni positivi non c'è nessuna variazione nei tempi di decadimento anche se, in dipendenza dal materiale attraversato, il muone potrebbe legarsi con un elettrone e creare lo stato di muonio ma questo processo ha un branching ratio completamente trascurabile.
 